
%includo il file che contiene la versione dei documenti
\newcommand{\versioneAnalisiDeiRequisiti}{3.2.0}			
\newcommand{\versioneNormeDiProgetto}{3.2.0}			
\newcommand{\versioneGlossario}{3.2.0}			
\newcommand{\versionePianoDiQualifica}{3.2.0}			
\newcommand{\versionePianoDiProgetto}{3.2.0}	
\newcommand{\versioneStudioDiFattibilita}{1.2.0}
\newcommand{\versioneSpecificaTecnica}{3.2.0}


\newcommand{\Versione}{\versioneSpecificaTecnica{}}	  %Versione Finale
\newcommand{\Data}{xxxx-xx-xx}				           %Data di creazione
\newcommand{\DataUltimaModifica}{xxxx-xx-xx}
\newcommand{\TipoDocumento}{Specifica Tecnica}	       %tipo documento

%includo il file header.tex (logo grande in prima pagina piu qualche altra regola)
%questo file contiene impostazioni comuni per tutte i documenti

%definizione packages utilizzati
\documentclass[a4paper]{article}
\usepackage[utf8x]{inputenc}
\usepackage{enumitem}
\usepackage[italian]{babel}
\usepackage{latexsym}
\usepackage{xparse}
\usepackage{float}
\usepackage{subfloat}
\usepackage{subfig}
\usepackage{fancyhdr}
\usepackage{eurofont}
\usepackage{lastpage}
\usepackage{graphicx}
\usepackage{textcomp}
\usepackage{booktabs}
\usepackage{color}
\usepackage{lscape}
\usepackage{hyperref}
%\usepackage[bookmarks=false]{hyperref}
\hypersetup{colorlinks=true, linkcolor=black, anchorcolor=red, urlcolor=blue}
\usepackage{longtable}
\usepackage{tabularx}
\usepackage{abstract}
\usepackage{appendix}
\usepackage{multicol}
\usepackage{bmpsize}
\usepackage[all]{hypcap}
\usepackage{titlesec}
\usepackage{indentfirst}
\usepackage{lipsum,titletoc}

%\setcounter{secnumdepth}{4}

%****************INIZIO GESTIONE SUBSECTION MULTIPLE
\makeatletter
\newcommand\level[1]{%
\ifcase#1\relax\expandafter\chapter\or
\expandafter\section\or
\expandafter\subsection\or
\expandafter\subsubsection\else
\def\next{\@level{#1}}\expandafter\next
\fi}
\newcommand{\@level}[1]{%
\@startsection{level#1}
{#1}
{\z@}%
{-3.25ex\@plus -1ex \@minus -.2ex}%
{1.5ex \@plus .2ex}%
{\normalfont\normalsize\bfseries}}

\newdimen\@leveldim
\newdimen\@dotsdim
{\normalfont\normalsize
\sbox\z@{0}\global\@leveldim=\wd\z@
\sbox\z@{.}\global\@dotsdim=\wd\z@
}

\newcounter{level4}[subsubsection]
\@namedef{thelevel4}{\thesubsubsection.\arabic{level4}}
\@namedef{level4mark}#1{}
\def\l@section{\@dottedtocline{1}{0pt}{\dimexpr\@leveldim*4+\@dotsdim*1+6pt\relax}}
\def\l@subsection{\@dottedtocline{2}{0pt}{\dimexpr\@leveldim*5+\@dotsdim*2+6pt\relax}}
\def\l@subsubsection{\@dottedtocline{3}{0pt}{\dimexpr\@leveldim*6+\@dotsdim*3+6pt\relax}}
\@namedef{l@level4}{\@dottedtocline{4}{0pt}{\dimexpr\@leveldim*7+\@dotsdim*4+6pt\relax}}

\count@=4
\def\@ncp#1{\number\numexpr\count@+#1\relax}
\loop\ifnum\count@<100
\begingroup\edef\x{\endgroup
\noexpand\newcounter{level\@ncp{1}}[level\number\count@]
\noexpand\@namedef{thelevel\@ncp{1}}{%
\noexpand\@nameuse{thelevel\@ncp{0}}.\noexpand\arabic{level\@ncp{1}}}
\noexpand\@namedef{level\@ncp{1}mark}####1{}%
\noexpand\@namedef{l@level\@ncp{1}}%
{\noexpand\@dottedtocline{\@ncp{1}}{0pt}{\the\dimexpr\@leveldim*\@ncp{5}+\@dotsdim*\@ncp{0}\relax}}}%
\x
\advance\count@\@ne
\repeat
\makeatother
\setcounter{secnumdepth}{100}
\setcounter{tocdepth}{100}
%****************FINE GESTIONE SUBSECTION MULTIPLE

%impostazioni relative alla visualizzazione delle section
%nell'indice
\titlecontents{section}
[0pt]%left indent
{\bfseries}
{\contentslabel{2.3em}}
{\hspace*{-2.3em}}
{\hfill\contentspage}
[]%separator


\oddsidemargin=.15in
\evensidemargin=.15in
\textwidth=6in
\topmargin=-.5in
\parindent=0in
\headheight=1in
\DeclareMathSizes{10}{10}{10}{10} %per piano qualifica
\pagestyle{fancy}
\lhead{
\bfseries {\Large \TipoDocumento}\\
\bfseries Versione: \Versione\\
}
\chead{}
\lhead{
\includegraphics[scale=0.455]{../Logo&Header/apertureHead.png}
}
\lfoot{\bfseries \TipoDocumento{} v\Versione}
\cfoot{}
\rfoot{\thepage\ of \mypageref{LastPage}}
\newcommand{\mypageref}[1]{
\hypersetup{linkcolor=black}\pageref{#1}\hypersetup{linkcolor=black}}
%\userpackage{lipsum}
\renewcommand{\footrulewidth}{0.4pt}

%definizioni comandi comuni utilizzati
\newcommand{\numref}[1]{\textsl{\nameref{#1} (\ref{#1})}}
\newcommand{\NomeGruppo}{Aperture Software}
\newcommand{\Progetto}{MaaP: MongoDB as an admin Platform}
\newcommand{\Prop}{CoffeeStrap}

%definizione tecnologie
\newcommand{\Node}{Node.js}
\newcommand{\NodeJS}{Node.js}
\newcommand{\Nodejs}{Node.js}

\newcommand{\mongodb}{MongoDB}

%tanti sub quanti ne vogliamo! :)
\newcommand{\subsubsubsection}{\level{4}}
\newcommand{\subsubsubsubsection}{\level{5}}
\newcommand{\subsubsubsubsubsection}{\level{6}}
\newcommand{\subsubsubsubsubsubsection}{\level{7}}
\newcommand{\subsubsubsubsubsubsubsection}{\level{8}}


%definizione comando per parola glossario
\newcommand{\gloss}[1]{\emph{#1}\ped{\emph{\tiny{G}}}}

\newcommand{\grassetto}{\textbf}

%per inserire immagini
\newcommand{\immagine}[2]{
\begin{center}
\begin{figure}[H]
\includegraphics[width=\textwidth]{{{#1}}}
\caption{#2}
\label{#1}
\end{figure}
\end{center}
}

\newcommand{\Glossario}{
Al fine di evitare ogni ambiguità nella comprensione del linguaggio utilizzato nel presente documento e, in generale, nella documentazione fornita dal gruppo \NomeGruppo{}, ogni termine tecnico, di difficile comprensione o di necessario approfondimento verrà inserito nel documento \emph{Glossario\_{}v\versioneGlossario{}.pdf}.\\
Saranno in esso definiti e descritti tutti i termini in corsivo e allo stesso tempo marcati da una lettera "G" maiuscola in pedice nella documentazione fornita.
}

\newcommand{\Prodotto}{
Lo scopo del prodotto è produrre un \gloss{framework} per generare interfacce web di amministrazione dei dati di business basato sullo stack \gloss{\Nodejs{}} e \gloss{\mongodb{}}.\\
L'obiettivo è quello di semplificare il lavoro allo \gloss{sviluppatore} che dovrà rispondere in modo rapido e standard alle richieste degli esperti di \gloss{business}.
}

%inizio pagina del documento
\begin{document}
\thispagestyle{empty}

\begin{center}\centerline{
%inserisco il logo grande della prima pagina
\includegraphics[scale=0.8]{../Logo&Header/logo.png}}

%metto il link dell'email sotto al logo
%{\href{mailto:ApertureSWE@gmail.com}{\color[rgb]{0.39,0.37,0.38}%ApertureSWE@gmail.com}}\\ [3pc]

\vspace{0.5in}

%titolo del progetto
{\Huge {\Progetto}}\\[.5pc]

\underline{\hspace{6in}}\\[8pc]

{\Huge {\TipoDocumento}}\\[1pc]
%{\emph{Versione \Versione}}\\
\end{center}

%\vspace{.05in}


%\vspace{.05in}

%informazioni documento
\begin{center}
%\section{Informazioni documento}
\begin{tabular}{r|l}
%\textbf{Nome} &\TipoDocumento \\
\textbf{Versione} & \Versione{} \\
\textbf{Data creazione} & \Data{2014-01-24} \\
\textbf{Data ultima modifica} & \DataUltimaModifica{} \\
\textbf{Stato del Documento} & Formale \\		          %CAMBIARE QUI
\textbf{Uso del Documento} & Esterno \\			          %CAMBIARE QUI
\textbf{Redazione} &  Pinato Giacomo,nome2,...\\		        %CAMBIARE QUI
\textbf{Verifica} & nome1,nome2,...\\			        %ED ANCHE QUI!
\textbf{Approvazione} & nome1,nome2,...\\				 %CAMBIARE QUI
\textbf{Distribuzione} & \parbox[t]{4cm}{\NomeGruppo{}}\\
\end{tabular}
\end{center}

\vspace{0.05in}

%inizio sommario del documento
\begin{abstract}
\begin{center}
Questo documento si propone di presentare la Specifica tecnica e architetturale per la Realizzazione del prodotto \textbf{MaaP}.
\end{center}
\end{abstract}

%\vspace{.4in}

%seconda pagina, diario delle modifiche
\newpage
Diario delle modifiche
\begin{center}
\begin{longtable}{|c|c|c|p{0.5\linewidth}|}
\toprule
\textbf{Versione} & \textbf{Data} & \textbf{Autore} & \textbf{Modifiche effettuate}\\

%aggiungere qui una midrule per aggiungere una nuova riga alla tabella
\midrule
1.2.0 & 2014-02-xx & ... (RE) & Approvazione documento\\
\midrule
1.1.0 & 2014-02-xx & ... (VE) & Verifica documento\\
\midrule
1.0.0 & 2014-03-05 & Giacomo Pinato (PR) & Componenti e Classi\\
\midrule
1.0.0 & 2014-01-25 & Giacomo Pinato (PR) & Tecnologie Utilizzate\\
\midrule
1.0.0 & 2014-01-25 & Fabio Miotto (PR) & Tecnologie Utilizzate\\
\midrule
1.0.0 & 2014-01-24 & Giacomo Pinato (PR) & Prima stesura del documento\\

\bottomrule
\caption{Registro delle modifiche}
\label{tab:changelog}

\end{longtable}
\end{center}

%terza pagina Indice (viene aggiornato in automatico con due compilazioni)
\newpage
\tableofcontents

%pagine successive hanno la lista di tabelle e lista delle figure
%(vengono aggiornate in automatico)
%\newpage
%\listoftables
%\listoffigures

%qui inizia la prima pagina ufficiale
\newpage
\section{Introduzione}
\label{1.0}
\subsection{Scopo del documento}
\label{1.1}
Il presente documento ha lo scopo di definire la progettazione ad alto livello del progetto \textbf{MaaP}, a partire dai requisiti individuati durante l'Analisi. Verrà presentata l'architettura generale secondo la quale saranno organizzate le varie componenti software, i \gloss{Design Pattern} e le tecnologie utilizzate per poi descrivere più dettagliatamente le varie componenti e relative dipendenze.

\subsection{Scopo del prodotto}
\label{1.2}
\Prodotto{}

\subsection{Glossario}
\subsubsection{Glossario2}
\subsubsubsection{Glossario3}
\subsubsubsubsection{Glossario4}
\label{1.3}
\Glossario{}

\subsection{Riferimenti}
\label{1.4}

\subsubsection{Normativi}
\label{1.4.1}
\begin{itemize}
\item \grassetto{Analisi dei requisiti}: Analisi\_{}dei\_{}Requisiti\_{}v\versioneNormeDiProgetto{}.pdf
\item Norme di Progetto: Norme\_{}di\_{}Progetto\_{}v\versioneNormeDiProgetto{}.pdf  (allegato alla presente documentazione)\\
\end{itemize}

\subsubsection{Informativi}
\label{1.4.2}
\begin{itemize}
\item Learning Node: O'Reilly Shelley Powers
\item AngularJS: O'Reilly Brad Green e Shyam Seshadri
\item Software Engineering (8th edition), Ian Sommerville, Pearson Education | Addison-Wesley
\item Design Patterns, E. Gamma, R. Helm, R. Johnson, J. Vlissides, Pearson Education | Addison-Wesley
\item Dall'idea al codice con UML 2       L. Baresi, L. Lavazza, M. Pianciamore, Pearson Education
\end{itemize}

\newpage
\section{Tecnologie utilizzate}
In questa sezione verrano elencate e descritte le tecnologie che si utilizzeranno durante lo sviluppo del progetto. In particolare la colonna portante del progetto sar\`{a} lo stack MEAN, ovvero MongoDB, Express, Angular e Node.js.

\subsection{MongoDB}
Il database con il quale la nostra applicazione dovrà interagire è realizzato con MongoDB, come specificato nel capitolato. Questa tecnologia offre i seguenti vantaggi:
\begin{itemize}
\item Facile indicizzazione: Ogni campo in MongoDB puo' diventare un indice;
\item Bilanciamento di carico: MongoDB scala orizzontalmente molto facilmente grazie all'utilizzo di shard;
\item Integrazione con Javascript: Query o altre funzioni scritte in Javascript possono essere eseguite direttamente dal dabatase;
\end{itemize}

\subsection{Javascript}
Si è deciso di utilizzare Javascript in quanto è il linguaggio su cui si basano tutte le altre tecnologie che andremo ad utilizzare, e offre quindi una facile integrazione, oltre ad essere un ottimo linguaggio per applicazioni web e client side.


\subsection{NodeJs}
Si è deciso di utilizzare il linguaggio Node.js in quanto consigliato dal capitolato e adatto al progetto. Le sue caratteristiche piu' vantaggiose sono:
\begin{itemize}
\item Modello event-driven: ovvero "programmazione ad eventi", che si basa su un concetto semplice: il flusso del programma non segue un corso specifico ma è guidato dalle azioni dell'utilizzatore;
\item Modello asincrono: grazie a questa caratteristica è possibile ridurre al minimo i tempi di morti in quanto, nell’attesa dell completamento di una operazione, si procede con altri flussi logici. 
\item Grande scalabilità: Grazie al modo in cui è implementato, Node.js riesce ad essere largamente scalabile con minimo sforzo.
\end{itemize}

\subsection{jQuery}
Per migliorare e semplificare la scrittura di funzioni nel linguaggio JavaScript, si è deciso di adottare il framework jQuery. I vantaggi sono:
\begin{itemize}
\item Semplicità: l'utilizzo di jQuery semplifica e facilita la scrittura di codice JavaScript, inoltre offre plug-in on-line per fornire nuove funzionalità.
\end{itemize}
Svantaggi:
\begin{itemize}
\item Pericolosità: jQuery offre funzionalità e plug-in molto utili, ma non tutto è compatibile con i vari browser, inoltre alcune funzionalità possono essere vecchie, non aggiornate o scritte male.
\end{itemize}

\subsection{JSON}
Rappresenta il tipo di messaggi con cui client e server si scambiano informazioni. I vantaggi offerti sono:
\begin{itemize}
\item Semplicità: i messaggi JSON sono più corti rispetto ad altri formati di interscambio, e vengono eseguiti più velocemente dal parser. JSON inoltre risulta più semplice e immediato rispetto ad esempio a XML.
\end{itemize}
Svantaggi:
\begin{itemize}
\item Restrittività: JSON è meno restrittivo rispetto ad XML, e questo può permettere di inserire errori nello scambio di messaggi.
\end{itemize}

\subsection{AngularJs}

\begin{itemize}
\item Two Way Data-Binding: Una delle caratteristiche principali di angular. Le modifiche apportate al model si rifletto direttamente sugli elementi del DOM, e le modifiche al DOM si ripercuotono automaticamente sul model. Questo alleggerisce tremendamente il codice necessario a controllare ad ascoltare e gestire il DOM, automatizzando il processo. E noi sappiamo che automatico \`{e} bene.
\item Templates: I template HTML sono parsati dal browswe nel DOM,il quale costituisce poi l’input per il compilatore Angular. Quest’ultimo poi crea il data binding tra il DOM e lo scope dei dati. Uno dei piu’ grandi vantaggi di questa tecnica e’ che separa presentazione da implementazione, in quanto i template html possono modificati senza alterare il modo in cui sono inseriti i dati.
\item Dependecy Injection: Angular possiede nativamente una “dependecy injection”, che aiuta gli sviluppatori facilitando la creazione, la comprensione e il testing dell’applicazione.
\item Directives: Le directives possono essere usate per definire tag HTML personalizzati che fungono da widget. Possono inoltre essere usate per “decorare” elementi con comportamenti personalizzati o per manipolare attributi del DOM. 

\end{itemize}

\subsection{HTML5}

\newpage
\section{Descrizione architettura}
\subsection{Metodo e formalismo di specifica}
Si è deciso di procedere utilizzando un approccio top-down per l’esposizione dell’architettura dell’applicazione, ovvero descrivendo inizialmente le componenti in generale per poi arrivare a trattarle al particolare.
Si descriveranno i package e i componenti per poi dettagliare le singole classi, specificando per ciascuna di esse il tipo, l’obiettivo e la funzionalità. Poi si passerà ad illustrare degli esempi d’uso di Design Pattern (descritti approfonditamente nell’ Appendice A) e le tecnologie utilizzate.


\subsection{Architettura generale}
-	L'architettura del framework segue un modello di architettura in stile three-tier che prevede la suddivisione dell'applicazione  in tre diversi strati dedicati rispettivamente all'interfaccia utente (Client), alla business logic (Controller) e alla gestione dei dati persistenti (Model). La parte Client segue il design pattern MVVM utilizzato da AngularJS ed è quindi suddivisa in Model, View, ViewModel.\\


Il seguente diagramma rappresenta l'architettura ad alto livello del framework, indicando i package e le relazioni che intercorrono tra questi.

...TODO aggiungere diagramma dell'installer CLI comprensivo di descrizione...

\immagine{./Diagrammi/GeneralePackage}{Architettura generale del software - vista package}
Nel precedente diagramma sono presenti le relazioni tra i package Client, Controller e ModelServer.\\
Vengono inoltre presentati tutti i sotto-package così da facilitare la comprensione dell'intero sistema.
\immagine{./Diagrammi/GeneraleClassi}{Architettura generale del software}
Nel precedente diagramma è presente l'architettura ad alto livello del software e vengono indicate le classi fondamentali per rappresentare le relazioni del modello  three-tier. TO DO DESCRIVERE 3TIER IN APPENDICE. I diagrammi di sequenza relativi allo scambio di segnali, lo scopo ed il contesto di utilizzo sono presenti nella sezione ????.

\newpage
\subsubsection{ModelServer}
\immagine{./Diagrammi/ModelClassi}{Diagramma delle classi del ModelServer}
Nel ModelServer sono presenti oggetti che rappresentano:
\begin{itemize}
\item Il database di analisi e quello degli utenti;
\item La gestione del file DSL e il suo parsing;
\item La gestione dei dati richiesti dal controller.
\end{itemize}
Tutte le operazioni di gestione, modifica e recupero dei dati vengono messe a disposizione dal model. In tal modo il controller è responsabile solamente di gestire la logica dell’applicazione.

\subsubsection{Controller}
\immagine{./Diagrammi/ControllerClassi}{Diagramma delle classi del Controller}
Il controller è responsabile dell’autenticazione delle richieste e del loro routing da Client a ModelServer e viceversa.

\subsubsection{Client}
\immagine{./Diagrammi/ClientClassi}{Diagramma delle classi del Client}
Nel Client sono presenti oggetti che rappresentano:
\begin{itemize}
\item I template per le pagine web;
\item I Controller per la gestione dei template;
\item Lo Scope per l'aggiornamento dei dati dei template;
\item I Servizi utilizzati dai Controller.
\end{itemize}

\newpage
\section{Componenti e Classi}

\subsection{MaaP}
\subsubsection{Informazioni sul package}
\immagine{./Diagrammi/ComponentiMaaP}{Componenti MaaP}
\paragraph{Descrizione}
Namespace globale per il progetto. Le relazioni tra i package ModelServer, Controller e Client identificano il modello di architettura three-tier.

\paragraph{Sotto-componenti}
\begin{itemize}
\item MaaP::ModelServer
\item MaaP::Controller
\item MaaP::Client
\end{itemize}

\subsection{MaaP::ModelServer}
\subsubsection{Informazioni sul package}
\immagine{./Diagrammi/ComponentiModelServer}{Componente MaaP::ModelServer}
\paragraph{Descrizione}
Package per il componente ModelServer del modello di architettura three-tier.

\paragraph{Sottocomponenti}
\begin{itemize}
\item MaaP::ModelServer::DataManager;
\item MaaP::ModelServer::Database;
\item MaaP::ModelServer::DSL.
\end{itemize}

\subsubsection{MaaP::ModelServer::DataManager}
\paragraph{Informazioni sul package}
\immagine{./Diagrammi/ComponentiDataManager}{Componente MaaP::ModelServer::DataManager}
\paragraph{Descrizione}
Componente parte del ServerModel per la gestione dei dati.
\paragraph{Sotto-componenti}
\begin{itemize}
\item MaaP::ModelServer::DataManager::DataBaseAnalysisManager;
\item MaaP::ModelServer::DataManager::DatabaseUserManager;
\item MaaP::ModelServer::DataManager::IndexManager.
\end{itemize}
\paragraph{Classi}
\subparagraph{MaaP::ModelServer::DataManager::JSonComposer}
\grassetto{Descrizione}\\
Classe che costruisce un file JSON a partire dalla struttura di una Collection, o di un Document, e dai suoi dati.\\
\grassetto{Utilizzo}
Viene utilizzata dai DatabaseManager per costruire il file JSON da inviare al Controller.

\subparagraph{MaaP::ModelServer::DataManager::IDatabaseManager}
\grassetto{Descrizione}\\
Interfaccia che rappresenta il gestore dei database. Contiene tutte le operazioni che si possono effettuare sul database e l'elaborazione dei dati recuperati da essi.\\
\grassetto{Utilizzo}
Viene utilizzata per la gestione delle richieste inoltrate dal Controller.

\subparagraph{MaaP::ModelServer::DataManager::IDataRetriever}
\grassetto{Descrizione}\\
Interfaccia attraverso cui i DatabaseManager dialogano con i batabase. Contiene le operazioni di lettura e scrittura nei database.\\
\grassetto{Utilizzo}
Viene utilizzata per recuperare e inserire dati, sui database, su richiesta dei DataManager.

\subsubsection{MaaP::ModelServer::DataManager::DatabaseAnalysisManager}
\paragraph{Informazioni sul package}
\immagine{./Diagrammi/ComponentiDatabaseAnalysisManager}{Componente MaaP::ModelServer::DataManager::DatabaseAnalysisManager}
\paragraph{Descrizione}
Componente parte del DataManager per la gestione dei dati del database di analisi.
\paragraph{Classi}
\subparagraph{MaaP::ModelServer::DataManager::DatabaseAnalysisManager::DatabaseAnalysisManager}
\grassetto{Descrizione}\\
Classe che rappresenta il gestore dei database di analisi. Contiene tutte le operazioni che si possono effettuare sul database di analisi e l'elaborazione dei dati recuperati da essi.\\
\grassetto{Utilizzo}
Viene utilizzata per la gestione delle richieste, relative al database di analisi, inoltrate dal Controller.
\grassetto{Classi da cui eredita}
\begin{itemize}
\item MaaP::ModelServer::DataManager::IDatabaseManager;
\end{itemize}
\grassetto{Relazioni con altre classi}
     \grassetto{MaaP::ModelServer::DataManager::JSonComposer}
     Relazione uscente, utilizza un riferimento a un oggetto di tipo JsonComposer per ottenere il JSON da spedire. 
     \grassetto{MaaP::ModelServer::DSL::CollectionData}
     Relazione uscente, utilizza un riferimento a un oggetto CollectionData che contiene la struttura di un file di descrizione 




\paragraph{Classi}
\begin{itemize}
\item MaaP::Server::Controller::Dispatcher
\item MaaP::Server::Controller::FrontController
\item MaaP::Server::Controller::PassportAdapter
\item MaaP::Server::Controller::Passport (?)
\end{itemize}

\paragraph{Iterazioni con altri componenti}
\begin{itemize}
\item MaaP::Server::Model
\item MaaP::Client::Controller
\end{itemize}

\subsection{Model}
\subsubsection{Informazioni sul package}
\immagine{./Diagrammi/ComponentiModelServer}{Componenti Model}
\paragraph{Descrizione}
Package del model server-side. Il model contiene i database di analisi e degli utenti, il parser e i file di descrizione DSL e ulteriori file di configurazione, oltre a tutta la logica di recupero e gestione dei dati dei database. 

\paragraph{Sottocomponenti}
\begin{itemize}
\item MaaP::Server::Model::DataManager
\item MaaP::Server::Model::Database
\item MaaP::Server::Model::DSL
\end{itemize}

\subsection{DSL}
\subsubsection{Informazioni sul package}
\immagine{./Diagrammi/ComponentiDSL}{Componenti DSL}
\paragraph{Descrizione}
Package della gestione DSL. Contiene i file di descrizione per le collection e il parser per generare la loro struttura partendo dagli stessi.

\paragraph{Classi}
\begin{itemize}
\item MaaP::Server::Model::DSL::DSLParser
\item MaaP::Server::Model::DSL::DSLManager
\item MaaP::Server::Model::DSL::CollectionData
\item MaaP::Server::Model::DSL::DSLDescriptionFile
\item MaaP::Server::Model::DSL::ParserInterface (?)
\end{itemize}

\subsection{Database}
\subsubsection{Informazioni sul package}
\immagine{./Diagrammi/ComponentiDatabase}{Componenti Database}
\paragraph{Descrizione}
Package dei database presenti nel model. Contiene il database degli utenti, il database di analisi e la raccolta delle query effettuate pi frequentemente.\`{u}

\paragraph{Classi}
\begin{itemize}
\item MaaP::Server::Model::Database::MongooseDBAnalysis
\item MaaP::Server::Model::Database::DBAnalysis
\item MaaP::Server::Model::Database::Mongoose (?)
\item MaaP::Server::Model::Database::MongooseDBFramework
\item MaaP::Server::Model::Database::DBFramework
\item MaaP::Server::Model::Database::USer
\item MaaP::Server::Model::Database::Query
\end{itemize}

\subsection{DataManager}
\subsubsection{Informazioni sul package}
\immagine{./Diagrammi/ComponentiDataManager}{Componenti DataManager}
\paragraph{Descrizione}
Package del DataManager. Contiene la logica per la gestione dei database degli utenti e di analisi, per la gestione degli indici e per la composizione dei JSON. 

\paragraph{Sottocomponenti}
\begin{itemize}
\item MaaP::Server::Model::DataManager::DatabaseAnalysisManager
\item MaaP::Server::Model::DataManager::DatabaseUserManager
\item MaaP::Server::Model::DataManager::IndexManager
\end{itemize}

\paragraph{Classi}
\begin{itemize}
\item MaaP::Server::Model::DataManager::JSonComposer
\item MaaP::Server::Model::DataManager::IDatabaseManager
\item MaaP::Server::Model::DataManager::IDataRetrievers
\end{itemize}

\subsection{DatabaseAnalysisManager}
\subsubsection{Informazioni sul package}
\immagine{./Diagrammi/ComponentiDatabaseAnalysisManager}{Componenti DatabaseAnalysisManager}
\paragraph{Descrizione}
Package del gestore del database di analisi.

\paragraph{Classi}
\begin{itemize}
\item MaaP::Server::Model::DataManager::DatabaseAnalysisManager::DatabaseAnalysisManager
\item MaaP::Server::Model::DataManager::DatabaseAnalysisManager::DataRetrieverAnalysis
\end{itemize}

\subsection{DatabaseUserManager}
\subsubsection{Informazioni sul package}
\immagine{./Diagrammi/ComponentiDatabaseUserManager}{Componenti DatabaseUserManager}
\paragraph{Descrizione}
Package del gestore del database degli utenti.

\paragraph{Classi}
\begin{itemize}
\item MaaP::Server::Model::DataManager::DatabaseUserManager::DatabaseUserManager
\item MaaP::Server::Model::DataManager::DatabaseUserManager::DataRetrieverUser
\end{itemize}

\subsection{IndexManager}
\subsubsection{Informazioni sul package}
\immagine{./Diagrammi/ComponentiIndexManager}{Componenti IndexManager}
\paragraph{Descrizione}
Package del gestore degli indici.

\paragraph{Classi}
\begin{itemize}
\item MaaP::Server::Model::DataManager::IndexManager::IndexManager
\end{itemize}

\subsection{Client}
\subsubsection{Informazioni sul package}
\immagine{./Diagrammi/ComponentiClient}{Componenti Client}
\paragraph{Descrizione}
Package contente la componente client-side del progetto.

\paragraph{Sottocomponenti}
\begin{itemize}
\item MaaP::Client::View
\item MaaP::Client::ControllerModelView
\item MaaP::Client::ModelClient
\end{itemize}


\subsection{View}
\subsubsection{Informazioni sul package}
\immagine{./Diagrammi/ComponentiView}{Componenti View}
\paragraph{Descrizione}
Package contente la componente i template delle pagine html generate dal framework.

\paragraph{Sottocomponenti}
\begin{itemize}
\item MaaP::Client::View::Template
\end{itemize}


\subsection{Template}
\subsubsection{Informazioni sul package}
\immagine{./Diagrammi/ComponentiTemplate}{Componenti Template}
\paragraph{Descrizione}
Package contente la componente i template delle pagine html generate dal framework.

\paragraph{Classi}
\begin{itemize}
\item MaaP::Client::View::Template::SignIn
\item MaaP::Client::View::Template::SignUp
\item MaaP::Client::View::Template::AdminMainPageCollection
\item MaaP::Client::View::Template::UserMainPageCollection
\item MaaP::Client::View::Template::MainPageDocument
\item MaaP::Client::View::Template::MainPageDocumentEdit
\item MaaP::Client::View::Template::UserProfileEdit
\item MaaP::Client::View::Template::UserProfile
\item MaaP::Client::View::Template::AdminProfile
\item MaaP::Client::View::Template::PasswordRecovery
\end{itemize}


\subsection{ControllerModelView}
\subsubsection{Informazioni sul package}
\immagine{./Diagrammi/ComponentiControllerModelView}{Componenti ControllerModelView}
\paragraph{Descrizione}
Package contente la componente i template delle pagine html generate dal framework.

\paragraph{Sottocomponenti}
\begin{itemize}
\item MaaP::Client::ControllerModelView::ControllerClient
\item MaaP::Client::ControllerModelView::Scope
\end{itemize}


\subsection{ControllerClient}
\subsubsection{Informazioni sul package}
\immagine{./Diagrammi/ComponentiControllerClient}{Componenti ControllerClient}
\paragraph{Descrizione}
Package contente la componente i template delle pagine html generate dal framework.

\paragraph{Classi}
\begin{itemize}
\item MaaP::Client::ControllerModelView::ControllerClient::ControllerAutenticazione
\item MaaP::Client::ControllerModelView::ControllerClient::ControllerCollection
\item MaaP::Client::ControllerModelView::ControllerClient::ControllerDocument
\item MaaP::Client::ControllerModelView::ControllerClient::ControllerProfilo
\item MaaP::Client::ControllerModelView::ControllerClient::ControllerMenu
\end{itemize}


\subsection{Scope}
\subsubsection{Informazioni sul package}
\immagine{./Diagrammi/ComponentiScope}{Componenti Scope}
\paragraph{Descrizione}
Package contente la componente i template delle pagine html generate dal framework.

\paragraph{Classi}
\begin{itemize}
\item MaaP::Client::ControllerModelView::Scope::Collection
\item MaaP::Client::ControllerModelView::Scope::Query
\item MaaP::Client::ControllerModelView::Scope::Document
\item MaaP::Client::ControllerModelView::Scope::Profilo
\item MaaP::Client::ControllerModelView::Scope::Menu
\end{itemize}


\subsection{ModelClient}
\subsubsection{Informazioni sul package}
\immagine{./Diagrammi/ComponentiModelClient}{Componenti ModelClient}
\paragraph{Descrizione}
Package contente la componente i template delle pagine html generate dal framework.

\paragraph{Sottocomponenti}
\begin{itemize}
\item MaaP::Client::ModelClient::Model
\item MaaP::Client::ModelClient::Services
\end{itemize}

\subsection{Services}
\subsubsection{Informazioni sul package}
\immagine{./Diagrammi/ComponentiServices}{Componenti Services}
\paragraph{Descrizione}
Package contente la componente i template delle pagine html generate dal framework.

\paragraph{Classi}
\begin{itemize}
\item MaaP::Client::ModelClient::Services::HTTP
\end{itemize}


\subsection{Model}
\subsubsection{Informazioni sul package}
\immagine{./Diagrammi/ComponentiModel}{Componenti Model}
\paragraph{Descrizione}
Package contente la componente i template delle pagine html generate dal framework.

\paragraph{Classi}
\begin{itemize}
\item MaaP::Client::ModelClient::Model::SessionData
\end{itemize}


\newpage
\section{Diagrammi di attività}

\newpage
\section{Design Pattern Utilizzati}%7.0
\subsection{Design Pattern architetturali} %7.1
\begin{itemize}
\item MVC: Utilizzato per l'architettura generale dell'applicazione.
\end{itemize}
\subsection{Design Pattern creazionali} %7.2
\begin{itemize}
\item Singleton
\end{itemize}
\subsection{Design Pattern comportamentali} %7.3
\begin{itemize}
\item Factory
\end{itemize}
\subsection{Design Pattern strutturali} %7.3
\begin{itemize}
\item Facade
\item Adapter
\end{itemize}

\newpage
\section{Stime di fattibilità e di bisogno di risorse}%8.0

\newpage
\section{Tracciamento} %9.0
\subsection{Tracciamento componenti - requisiti} %9.1
\subsection{Tracciamento requisiti - componenti} %9.2

\newpage
\appendix
\section{Descrizione Design Pattern} %A.0

\subsection{Design Pattern architetturali} %A.1
\subsection{Design Pattern 1} %A.1.1
\begin{itemize}
\item \grassetto{Scopo:} 
\item \grassetto{Motivazione:} 
\item \grassetto{Applicabilità:}
\end{itemize}

\subsection{Design Pattern creazionali} %A.2
\subsection{Design Pattern 1} %A.2.1
\begin{itemize}
\item \grassetto{Scopo:} 
\item \grassetto{Motivazione:} 
\item \grassetto{Applicabilità:}
\end{itemize}

\subsection{Design Pattern strutturali} %A.3
\subsection{Design Pattern 1} %A.3.1
\begin{itemize}
\item \grassetto{Scopo:} 
\item \grassetto{Motivazione:} 
\item \grassetto{Applicabilità:}
\end{itemize}

\subsection{Design Pattern comportamentali} %A.4
\subsection{Design Pattern 1} %A.4.1
\begin{itemize}
\item \grassetto{Scopo:} 
\item \grassetto{Motivazione:} 
\item \grassetto{Applicabilità:}
\end{itemize}

%FINE DOCUMENTO NON CANCELLARE
\end{document}
